\section{概率推断与统计保证：S14--S17}
\label{sec:statistical}

本节统计回归示例采用净注入为正的 $p,q$，故写成 $v=b\mathbf1+Rp+Xq$；它与总论中净负荷为正的约定相差一个符号。$v$ 表示所选电压观测坐标，幅值或平方电压的比例系数需相应纳入 $R,X$。本节区分四种问题：怎样把损失转为信念权重，怎样在正确统计模型下构造有限样本置信集合，怎样在树参数空间上进行后验计算，以及怎样校准未来观测的预测集合。它们分别对应 S14、S15、S16、S17，保证对象并不相同。以下凡标注“本文推导”的内容，均为针对末端 $P/Q/V$ 与隐藏零注入最简有根树所作的独立数学解释，不能归作被引文献已经证明的配电网结论。

\subsection{S14：Bissiri--Holmes--Walker，一般损失驱动的信念更新}
\label{subsec:S14}

\paragraph{文献定位与经核实的原文事实。}
\citet{S14} 的研究对象是由总体期望损失最小值定义的参数，允许不为全部观测指定完整似然。其第 1.1 节式 (7)--(8) 使用期望损失与相对先验的 KL 散度；第 3 节讨论两项的相对尺度。文中的一致更新指分批与合并更新相容。附录关于 KL 的必要性有其散度类和规则假设，不能表述为“任意理性推断都只能采用 KL”。以下编号依据 arXiv v2 全文，期刊书目信息见参考文献。

\paragraph{本文推导：从变分目标得到树的 Gibbs 权重。}
令 $\vartheta=(T,\theta_T)$ 包含最简树及其正边参数，先验为 $\pi$，观测为 $z_{1:n}$，预先固定的累积损失为
\[
 L_n(\vartheta)=\sum_{t=1}^n\ell(\vartheta;z_t),\qquad \eta>0.
\]
在所有 $\nu\ll\pi$ 的概率分布上考虑
\[
 J(\nu)=\eta\int L_n(\vartheta)\,\nu(d\vartheta)
       +\mathrm{KL}(\nu\Vert\pi).
\]
若 $0<Z_n:=\int e^{-\eta L_n}\,d\pi<\infty$，定义
\[
 \pi_n(d\vartheta)=Z_n^{-1}e^{-\eta L_n(\vartheta)}\pi(d\vartheta).
\]
将 $\log(d\nu/d\pi_n)=\log(d\nu/d\pi)+\eta L_n+\log Z_n$
代入即得
\[
 J(\nu)=\mathrm{KL}(\nu\Vert\pi_n)-\log Z_n.
\]
KL 的非负性给出唯一最优分布 $\nu=\pi_n$，其中唯一性按概率测度理解。
这段推导解释了指数权重来自什么优化问题，并没有证明 $\pi_n$ 的任意可信集合具有频率覆盖率。

若只保留有限候选树 $\mathcal T_{\mathrm{cand}}$，可对连续参数积分：
\[
 w_n(T)=
 \frac{\pi(T)\int e^{-\eta L_n(T,\theta)}\pi(d\theta\mid T)}
 {\sum_{S\in\mathcal T_{\mathrm{cand}}}\pi(S)
       \int e^{-\eta L_n(S,\xi)}\pi(d\xi\mid S)}.
\]
用 $\exp[-\eta\min_\theta L_n(T,\theta)]$ 代替积分，会丢掉不同树上可接受参数区域的体积；它是另一种评分规则，通常不等于上式。固定两点参数时，后验对数密度比满足
\[
 \log\frac{\pi_n(\vartheta_a)}{\pi_n(\vartheta_b)}
 =\log\frac{\pi(\vartheta_a)}{\pi(\vartheta_b)}
  -\eta\{L_n(\vartheta_a)-L_n(\vartheta_b)\},
\]
这里连续参数的密度须相对共同基准测度定义。因此温度和损失量纲都会直接改变权重集中程度。例如将电压由伏特改为千伏，未标准化的平方误差会缩小 $10^6$ 倍；保持相同数值 $\eta$ 就不是保持同一分析。

\paragraph{本文推导：顺序相容性需要保留什么。}
固定 $\eta$ 且损失可加时，两批数据 $A,B$ 的更新满足
\[
 e^{-\eta L_B}e^{-\eta L_A}\pi
 \ \propto\ e^{-\eta(L_A+L_B)}\pi.
\]
这一代数性质本身不要求各项统计独立，但不能由此获得依赖数据下的有效样本量或覆盖保证。若处理第二批时重新选择 $\eta$、改变残差标准差、使用前后不同的 RNJ 候选域，更新的目标已经变了。尤其把同一批数据得到的 cherry 先验再作为独立外部信息加入损失，会让同一证据被重复计权；必须记录先验来源及训练、评分的划分。

\paragraph{本文推导：可信区间的尺度为何不能自动校准。}
暂时固定拓扑并假设参数位于光滑内部，令总体损失具有唯一极小点 $\theta_\star$，记
\[
 H=\mathbb E[\nabla^2_\theta\ell(\theta_\star;Z)],\qquad
 J=\operatorname{Var}[\nabla_\theta\ell(\theta_\star;Z)].
\]
在独立同分布、可微、适当矩条件和局部近似成立时，经验损失极小点的渐近协方差为
$n^{-1}H^{-1}JH^{-1}$；而上述损失后验的局部二次近似协方差为
$(\eta nH)^{-1}$。两者相等要求
\[
 J=\eta^{-1}H.
\]
取 $H=I_2$、$J=\operatorname{diag}(1,9)$，便不存在同时校准两个方向的标量 $\eta$。时间相关性还会将 $J$ 换成长程协方差。这个计算只在已列局部条件下成立；它更不能被搬到内部边长为零、树结构发生改变的非光滑边界。

\paragraph{本文推导：对 Topo 的适当迁移。}
可定义统一量纲的电压残差损失，例如
\[
 \ell(T,\theta;z_t)
 =\rho\!\left(
 \left\|S_t^{-1/2}\{v_t-b_t\mathbf1-R_Tp_t-X_Tq_t\}\right\|_2
 \right),
\]
其中正定尺度矩阵 $S_t$、根电压项 $b_t$ 的处理规则以及功率符号约定都应预先说明。若 $S_t$ 由训练数据估计，评分时应冻结；否则需要完整说明自适应规则。平方损失、Huber 损失及不同时间权重分别定义不同的总体最优目标。只有该目标唯一对应目标最简树、真树具有先验支持且统计条件满足时，才有继续讨论一致性的基础。

特别地，若 $T_\star\notin\mathcal T_{\mathrm{cand}}$，则
$w_n(T_\star)=0$ 对所有 $n$ 恒成立。若两棵树在观测设计上满足
$R_{T_1}p_t+X_{T_1}q_t=R_{T_2}p_t+X_{T_2}q_t$，
它们的残差损失相同，后验差异只能来自先验和参数体积。更新规则不会补充缺失的激励或消除隐藏度二节点带来的观测等价。

\paragraph{复现审查点。}
应保存损失的完整定义、量纲、$\eta$、先验支持、候选生成数据、连续参数积分或近似方式，以及根电压的处理。至少比较温度变化、先验误报、候选漏真树及依赖增强四类敏感性。输出应称为“给定损失与候选域的信念权重”；只有另外完成覆盖率理论或独立校准，才进一步使用具有明确频率含义的置信语言。

\subsection{S15：Wasserman--Ramdas--Balakrishnan，Universal Inference}
\label{subsec:S15}

\paragraph{文献定位与经核实的原文事实。}
\citet{S15} 以拆分似然比构造有限样本置信集合与检验，避免依靠常规似然比的渐近卡方分布。核查的 arXiv v4 第 2 节 Theorem 1 与 Theorem 3 分别给出集合覆盖和检验有效性；第 6 节讨论条件似然，第 7 节讨论特定凸模型下的错设。“Universal” 不表示任意残差分数、错误似然和任意停止规则都自动有效。

\paragraph{本文推导：一个归一化积分是保证的核心。}
设真实分布为某个 $p_{\vartheta_\star}$，数据拆分为独立的训练部分 $D_1$ 和检验部分 $D_0$。利用 $D_1$ 构造任意归一化预测密度 $\widehat q_1$；它可以来自一个较差的估计器，甚至不同的备择模型。令
\[
 L_0(\vartheta)=\prod_{i\in D_0}p_\vartheta(Z_i),\qquad
 Q_1(D_0)=\prod_{i\in D_0}\widehat q_1(Z_i),\qquad
 E(\vartheta)=\frac{Q_1(D_0)}{L_0(\vartheta)}.
\]
假设检验样本在该模型下独立，条件于 $D_1$ 后 $\widehat q_1$ 固定，因此
\[
 \mathbb E_{\vartheta_\star}[E(\vartheta_\star)\mid D_1]
 =\int_{\{L_0(\vartheta_\star)>0\}}Q_1(z)\,dz\le1.
\]
由 Markov 不等式，
\[
 \mathbb P_{\vartheta_\star}\{E(\vartheta_\star)>1/\alpha\}\le\alpha,
 \qquad
 \mathcal C_\alpha=\{\vartheta:E(\vartheta)\le1/\alpha\}
\]
具有至少 $1-\alpha$ 的覆盖率。支持集不重合时积分可能严格小于一；不应将其无条件写成等号。$\exp(-\mathrm{SSE})$ 如果没有对应正确且归一化的观测模型，就不能直接放进这个积分证明。

\paragraph{本文推导：把置信集合投影到拓扑。}
树 $T$ 下的边参数、噪声参数等记为 $\theta\in\Theta_T$。定义
\[
 U_T=\frac{Q_1(D_0)}
 {\sup_{\theta\in\Theta_T}L_0(T,\theta)},\qquad
 \mathcal C^T_\alpha=\{T:U_T\le1/\alpha\}.
\]
若真实参数 $(T_\star,\theta_\star)$ 在所声明模型中，
分母至少为 $L_0(T_\star,\theta_\star)$，所以
$\mathbb E[U_{T_\star}]\le1$，进而
$\mathbb P(T_\star\in\mathcal C^T_\alpha)\ge1-\alpha$。
拓扑不可识别时不应将统计保证解读为唯一拓扑可恢复。观测等价的树可能同样难以排除；非空或单元素集合也必须结合模型与计算条件解读。

这一投影面对的是某一真实树的保留事件，不必仅因逐树反演而额外对全部候选树作 Bonferroni 校正。但是另行同时宣称大量边命题、或不断试验后挑选最有利的检验统计量，需要与该新报告对象匹配的联合有效性论证。

\paragraph{本文推导：计算证书的方向决定统计安全性。}
当精确求 $\sup L_0$ 很难时，若计算出经过验证的上界
\[
 B_T\ge\sup_{\theta\in\Theta_T}L_0(T,\theta),
\]
则 $\widetilde U_T=Q_1/B_T\le U_T$ 仍然安全，只是更保守。反之，局部优化得到的可行参数给出
$L_0(T,\widehat\theta)\le\sup L_0$，用它作分母会放大统计量，可能误删真树。

如果求解器处理负对数似然 $F_T(\theta)=-\log L_0(T,\theta)$，正确方向是
\[
 \underline F_T\le\inf_\theta F_T(\theta)
 \quad\Longrightarrow\quad
 e^{-\underline F_T}\ge\sup_\theta L_0(T,\theta).
\]
因此，在目标函数确为完整负对数似然、常数项及数值误差被处理的前提下，全域最小化的可验证下界可以进入安全分母；最小化的可行解上界通常不可以。RNJ 候选筛选后的最优值仅约束筛选域，不能充当全树域的似然上界。这个方向判断说明，MILP gap 的数值必须与具体统计分母对应，不能单独解释为拓扑覆盖证书。

若只检验了少数树，未检验树应标记为未排除，或将结论明确限制在有外部依据的模型域内。即使候选集合由独立训练集构造，独立性也不会让被漏掉的真树重新出现。候选域随机缩小和检验统计量可校准是两个不同的问题。

\paragraph{本文推导：末端 $P/Q/V$ 的可行模型与限制。}
在功率测量可当作已知设计量且残差模型成立时，可设
\[
 v_t\mid p_t,q_t
 \sim N\!\left(b_t\mathbf1+R_Tp_t+X_Tq_t,\Omega\right).
\]
为了应用拆分证明，需要正确处理根电压 $b_t$、未知 $\Omega$、功率测量误差和时序相关。训练集预测密度的均值或方差不能偷看待评分的 $v_t$ 后再调到最优。未知根电压若每个时刻都自由拟合，会损失辨识信息；只有分母 profile 化并不能为使用同一观测优化后的分子恢复归一化。若采用依赖模型，应使用归一化的
$q(D_0\mid D_1)/p_\vartheta(D_0\mid D_1)$，直接把相邻采样点看作独立样本一般缺少依据。

电压线性化有系统偏差时，真实分布可能不在树模型族中。有限个树高斯回归模型的并通常不构成密度意义下的凸集合，因此不能直接套用原文第 7 节的凸模型错设扩展。若为了凸化而允许树模型的混合，统计参数就成为混合分布；它又不必对应单一物理拓扑。

\paragraph{本文推导：多次拆分与持续监测。}
若 $E_1,\ldots,E_K$ 各自非负且真值处期望不超过一，则
$\overline E=K^{-1}\sum_k E_k$ 仍有该性质，彼此独立不是必要条件；最大值或观察结果后选择一个 $E_k$ 则没有这个简单保证。若要持续监测，可研究条件预测乘积
\[
 M_t(\vartheta)
 =\prod_{s=1}^t
 \frac{q_s(Z_s\mid Z_{<s})}{p_\vartheta(Z_s\mid Z_{<s})},
\]
其中 $q_s$ 必须在看到 $Z_s$ 前确定且归一化。此时才具备建立非负超鞅与停止时保证的基础，固定样本拆分公式本身不提供这些条件。

\paragraph{复现审查点。}
应保存训练与检验索引、完整似然、分子拟合来源、profile 参数域、求解上下界、数值容差，以及未检查拓扑的处理。验证应同时统计真树覆盖和集合大小，并设置弱激励、短内部边、重尾、时序相关及模型错设场景。覆盖很高但始终保留全域是有效而无辨识力的程序，不能只报告覆盖率。

\subsection{S16：Yao--Wu--Bharath--Baladandayuthapani，超度量协方差的树空间贝叶斯推断}
\label{subsec:S16}

\paragraph{文献定位与经核实的原文事实。}
\citet{S16} 将超度量协方差表示为树分裂与边长，构造树空间先验并进行高斯后验抽样。本文使用作者网站主文 PDF；arXiv v1 HTML 实际呈现补充材料。主文 Definition 1、Corollary 1 和第 5 节分别给出矩阵条件、分裂展开及高斯模型；演示先验对零内部边的边界不赋正质量。补充材料 S1.3 报告了所试重尾错设下可信区间覆盖变差。该经验现象不是一般错设定理。

\paragraph{本文推导：共同路径矩阵与树分裂表示。}
令 $A_e$ 是边 $e$ 下游可观测末端的集合，
$u_e\in\{0,1\}^m$ 为其指示向量。对每条边给正权重 $a_e$，则
\[
 K_T=\sum_{e\in E(T)}a_eu_eu_e^\top,\qquad
 (K_T)_{ij}=\sum_{e:\ i,j\in A_e}a_e.
\]
第二式正是根到末端 $i,j$ 最近公共祖先的累积权重。树分裂集合具有嵌套或不交的层级结构；任意选择若干二进制向量，不一定对应合法树。若每个末端有正的独有边长，
\[
 x^\top K_Tx=\sum_e a_e(u_e^\top x)^2
 \ge\sum_{i=1}^m a_{\{i\}}x_i^2>0
 \quad(x\ne0),
\]
故正定性直接成立。对三个末端，公共祖先层级给出
$(K_T)_{ij}\ge\min\{(K_T)_{ik},(K_T)_{jk}\}$。
这类超度量协方差允许对角元不同；它不要求所有末端到根等距。不要把矩阵中的相似度不等式和叶间距离的超度量不等式混淆。

在径向线性响应且支路电阻为正的模型中，$R_T$ 有同样形式，$X_T$ 在支路电抗满足相应正性条件时也有同样形式。共享的是分裂组合结构；电阻、电抗和随机增量方差有不同物理单位，不能直接共用边长先验。将无法区分的零注入隐藏度二节点压缩、保留可识别的隐藏分支节点，也不同于假设每条边具有独立高斯随机增量。

\paragraph{本文推导：原始电压协方差一般不是该树矩阵。}
若中心化电压近似 $v=Rp+Xq+\varepsilon$，则在噪声与功率不相关时
\[
 \operatorname{Cov}(v)
 =R\Sigma_pR^\top+X\Sigma_qX^\top
  +R\Sigma_{pq}X^\top+X\Sigma_{qp}R^\top+\Sigma_\varepsilon.
\]
即使 $R$ 是严格正边的共同路径矩阵，右端也不必超度量。取
\[
 R=\begin{pmatrix}3&2&1\\2&3&1\\1&1&2\end{pmatrix},
 \qquad
 \Sigma_p=\operatorname{diag}(1,2,1),\qquad q=\varepsilon=0.
\]
它对应共同根边长 $1$、末端 $1,2$ 的共同内部边长 $1$、三个独有末端边长均为 $1$。直接相乘得到
\[
 \operatorname{Cov}(v)
 =\begin{pmatrix}18&19&9\\19&23&10\\9&10&7\end{pmatrix}.
\]
由于 $9<\min(19,10)$，超度量不等式失效。故将实测 $v_t$ 直接当作论文的高斯树增量样本，不能仅以电网径向为由获得模型依据。较自然的迁移是把树空间表示用于 $R/X$，再为 $P/Q/V$ 建立条件观测模型，或为估计得到的响应矩阵建立包含估计误差的似然。

\paragraph{本文推导：后验计算的对象和边界。}
若观测确实为 $x_1,\ldots,x_n\stackrel{\mathrm{iid}}{\sim}N(0,K_T)$，
令 $S=n^{-1}\sum_tx_tx_t^\top$，其树参数对数后验除常数外为
\[
 -\frac n2\log\det K_T
 -\frac n2\operatorname{tr}(SK_T^{-1})
 +\log\pi(T)+\log\pi(a\mid T).
\]
固定拓扑时优化或抽样连续边长；改变拓扑时改变兼容分裂集合。一条内部边缩到零后可以跨到相邻分辨率的树域，这是几何化局部搜索的作用。但局部可达不等于有限时间已经充分遍历后验；若若干高质量模式隔着低后验区域，似然轨迹达到平台仍可能没有交换树模式。

还可以具体检查正边长更新。设单边先验为均值 $\lambda>0$ 的指数分布，用正半轴截断正态提议 $a'\mid a$，未截断中心为 $a$、标准差为 $\sigma$。若 $\Phi_N$ 是标准正态分布函数，则
\[
 \frac{q(a\mid a')}{q(a'\mid a)}
 =\frac{\Phi_N(a/\sigma)}{\Phi_N(a'/\sigma)}.
\]
因此，固定其余边和拓扑时，正确 MH 比率为
\[
 r=\frac{L(a')}{L(a)}
       \exp\!\left(-\frac{a'-a}{\lambda}\right)
       \frac{\Phi_N(a/\sigma)}{\Phi_N(a'/\sigma)}.
\]
虽然未截断高斯核关于 $a,a'$ 对称，截断后的归一化常数不对称；忽略最后一项会改变目标分布。这是本文对正边长提议的独立推导，也说明几何可行性与后验采样正确性要分别检查。

若所有内部边都服从连续正值先验，则
$\mathbb P(a_e=0\mid D)=0$，只要后验相对该先验绝对连续。于是，真实多分叉树在这种参数化下通常只能由很短的二叉分辨边逼近，不能由后验精确选中零边。针对最简非二叉树，需要显式为低维边界赋质量，或者清楚定义短边收缩阈值及其统计含义。根的不同出线若没有共同正边，响应矩阵也允许部分非对角元为零；把正共同根边作为模型条件前应核查物理根定义。

\paragraph{本文推导：为什么不能逐元素平均两棵树。}
令 $\mathbf1=(1,1,1)^\top$，
\[
 K_{12}=I+\mathbf1\mathbf1^\top+a\,u_{\{1,2\}}u_{\{1,2\}}^\top,\quad
 K_{13}=I+\mathbf1\mathbf1^\top+a\,u_{\{1,3\}}u_{\{1,3\}}^\top,\quad a>0.
\]
两者都为合法共同路径矩阵；它们算术平均的三个非对角元却是
$1+a/2,1+a/2,1$，违反超度量不等式。这说明一个预测协方差的后验均值可以有意义，却未必对应某棵物理树。若输出必须仍为树，就应定义树域内的损失最优代表，例如在正确指定的树距离下求 Fr\'echet 均值，并检查存在性、所用距离及边界。

\paragraph{所取主文版本的公式核验点与本文数学检查。}
作者网站 PDF 印刷第 7 页将扩展树距离写为
$d_{\mathrm{tree}}=d_{\mathrm{BHV}}+\|x-y\|_2$，
随后的 Theorem 1 声称相应空间为 CAT(0)。这已通过渲染 PDF 目视确认，不是网页公式提取造成的差异。按该加法原样理解，在固定树域内只改变一条内部边长 $a$ 和一条末端边长 $b$，距离就是 $|a-a'|+|b-b'|$。从 $(1,1)$ 到 $(2,2)$ 的对角线与先横后竖折线，都可参数化为长度 $2$ 的测地线，且中点不同。因此不能原样据此使用唯一测地线结论。一个应核查的修正方向是
$d_2=\sqrt{d_{\mathrm{BHV}}^2+\|x-y\|_2^2}$；
本文不将此修正冒充作者已经发布的勘误。

主文印刷第 15 页式 (9) 将 Metropolis--Hastings 接受概率写成
$\max\{1,r\}$。若目标密度为 $\pi_D$、提议为 $q$，正确的概率构造应为
\[
 a(T,T')=\min\!\left\{1,
 \frac{\pi_D(T')q(T\mid T')}{\pi_D(T)q(T'\mid T)}\right\}.
\]
此时 $\pi_D(T)q(T'\mid T)a(T,T')$ 等于正反向未接受概率流的最小值，因而对称并满足详细平衡；$\max\{1,r\}$ 则可大于一。以上只核查所下载版本的数学表达，不据此断言作者实现也采用错误概率，或否定分裂表示本身。源文件和关键页截图路径记录在随附来源说明中。

\paragraph{复现审查点。}
先验证输入究竟为树增量样本、响应矩阵估计，还是原始电压序列；再核实距离、接受概率、截断正态提议的正反向归一化因子、零边先验质量和输出树的定义。多个初始拓扑下应比较分裂访问频率、模式切换、有效样本量及后验预测，而非仅比较似然高度。候选 RNJ 树可提供初值，但若禁止离开候选域，得到的是截断后验；未遍历树的质量不能由这些样本估计。

\subsection{S17：Angelopoulos--Bates，共形预测及其拓扑迁移边界}
\label{subsec:S17}

\paragraph{文献定位与经核实的原文事实。}
\citet{S17} 是共形预测教程。核查的 arXiv v6 PDF 第 1.1 节和附录 D 给出 split conformal 构造及秩论证；Theorem D.1 给覆盖下界，Theorem D.2 的近似匹配上界另需连续性条件。第 3.1 节区分边际和条件覆盖，第 4 节讨论若干扩展。本文采用固定 v6 PDF 的 2022 年版本信息，不使用 HTML 页动态显示的日期作为发表时间。

\paragraph{本文推导：覆盖来自校准点和新点的秩对称性。}
在独立训练数据上冻结预测器及评分函数 $s(x,y)$。
有 $n$ 个校准样本 $(X_i,Y_i)$ 和一个新样本，并假设条件于训练数据后这些样本可交换。令
\[
 s_i=s(X_i,Y_i),\quad
 k=\lceil(n+1)(1-\alpha)\rceil,\quad
 \widehat q=
 \begin{cases}
  s_{(k)},&k\le n,\\
  +\infty,&k=n+1,
 \end{cases}
\]
其中 $s_{(k)}$ 是校准分数第 $k$ 个次序统计量，且 $0<\alpha<1$，构造
\[
 \mathcal C(X_{\mathrm{new}})
 =\{y:s(X_{\mathrm{new}},y)\le\widehat q\}.
\]
若分数几乎必然无并列，则新分数在全部 $n+1$ 个分数中的秩均匀分布，因而当 $k\le n$ 时
\[
 \mathbb P\{Y_{\mathrm{new}}\in\mathcal C(X_{\mathrm{new}})\}
 =\frac{k}{n+1}\ge1-\alpha.
\]
若存在并列，使用弱不等式的集合通常更保守，覆盖下界保留；不能不加条件继续声称同样精确的上界。若 $n=9$、$\alpha=0.05$，则 $k=10$，应取无限阈值而非把分位数参数硬截为 $1$ 后使用最大校准残差；后者一般只能获得 $9/10$ 的无并列覆盖率。

这个证明没有评价预测器是否正确。预测器较差时集合仍可通过变宽取得覆盖，因此覆盖率和集合大小必须一同报告。训练集可以复杂地拟合 RNJ 或树回归，关键是校准评分不能继续根据校准标签选择最有利的模型和阈值而又原样引用该证明。

\paragraph{本文推导：对未来电压的直接、较明确用法。}
把 $X=(p,q,\hbox{已知工况})$，$Y=v$。在训练集拟合
$\widehat v(X)$ 与正的尺度函数 $\widehat\sigma_j(X)$，可定义向量分数
\[
 s(X,v)=\max_{1\le j\le m}
 \frac{|v_j-\widehat v_j(X)|}{\widehat\sigma_j(X)}.
\]
共形集合于是为
\[
 \mathcal C_v(X)=
 \left\{v:\ |v_j-\widehat v_j(X)|
 \le\widehat q\,\widehat\sigma_j(X),\quad j=1,\ldots,m\right\}.
\]
它校准的是单个未来时刻整组末端电压同时落入区域的事件。若逐节点分别用 $\alpha$ 校准，则通常只得到每个节点自己的边际覆盖，不能自动宣称整组同时覆盖；最大分数是在评分层面直接定义了整组目标。若要覆盖一整段未来时间轨迹，可将整个窗口视为一个输出对象，再为窗口设置分数，但必须重新检查窗口之间的可交换性和有效样本数量。

根电压漂移、主动探测策略随历史改变、季节性负荷及滑动窗口重叠都可能破坏上述交换条件。简单随机打乱一条时间序列并不能创造交换性；基于块的经验改善也不是无需条件的精确有限样本定理。

\paragraph{本文推导：预测覆盖为什么不等于拓扑置信度。}
未来电压的集合事件是
$v_{\mathrm{new}}\in\mathcal C_v(X_{\mathrm{new}})$；
目标树的集合事件是
$T_\star\in\mathcal C_T(D)$。这两个随机对象不同。
若校准样本只是同一未知网络上不同时间的 $P/Q/V$，
每个时刻并没有一个可观测的独立拓扑标签。
将训练阶段得到的单棵 $\widehat T$ 固定，再为其电压残差作共形校准，不能检验 $\widehat T=T_\star$。

极端反例是取 $\mathcal C_v(X)=\mathbb R^m$，其电压覆盖恒为一，然而 $\widehat T$ 可以完全错误。更具物理意义的反例是末端功率设计只覆盖一个低维子空间；两棵不同树在这个子空间上产生相同电压。它们在这些工况上的预测均可良好，但预测校准并不打破拓扑不可识别性。

\paragraph{本文推导：如果确实要共形拓扑集合，需要什么数据。}
一种可明确陈述的路线是以整套网络实例为样本：
\[
 X_i=\hbox{网络 }i\hbox{ 的末端观测数据包},\qquad
 Y_i=T_i^\star,
\]
并拥有已知真树的、与目标网络同机制可交换的校准实例。对任何候选树 $T$ 定义评分 $s(X,T)$，例如冻结推断器提供的拓扑损失或负对数权重，然后构造
$\mathcal C_T(X)=\{T:s(X,T)\le\widehat q\}$。
此时才可由网络实例层面的秩对称性讨论拓扑边际覆盖。

若只在测试时把该集合与 RNJ 候选域相交，真树可能被额外删除，原保证不保留。即使校准时也使用候选生成器，也需要保证评分在整个标签空间有定义；为域外树赋 $+\infty$ 时，若候选遗漏频繁，正确校准可能要求 $\widehat q=+\infty$，此时必须按构造把全标签空间纳入。人工将输出硬截到候选域，会再次失去保证。

仿真网络标签可以帮助验证流程，但从仿真到真实网络存在模型差异时，交换性是需要检验或另行建模的假设。对某一目标网络的条件覆盖也比跨网络总体平均覆盖更强，不能仅从后者推导。

\paragraph{本文推导：边际覆盖允许局部系统失败。}
设常规工况概率为 $0.95$、异常工况概率为 $0.05$。若区域在常规工况总覆盖、在异常工况从不覆盖，总体覆盖仍为 $0.95$。这只是概率分解
\[
 \mathbb P(\mathrm{cover})
 =\sum_g\mathbb P(G=g)\mathbb P(\mathrm{cover}\mid G=g)
\]
的直接结果。所以针对电压越限附近、较弱激励、根电压异常或特定季节，应单独报告分组样本数和覆盖；不能把总体 $95\%$ 描述为每个工况都有 $95\%$。分组专门校准还会消耗各组校准样本，需重新应用对应的有限样本分位数规则。

\paragraph{复现审查点。}
应明确样本单位是时刻、窗口还是整网；冻结评分和尺度估计；保存准确的次序统计量索引、并列处理和无限阈值行为。报告预测区域大小、整体及分组覆盖、校准划分不确定性，并保持独立验证集。最终结果标签必须写清“未来电压预测覆盖”或“跨网络拓扑标签覆盖”，不能把任一形式简写成未经定义的拓扑可信度。

\paragraph{四篇方法的组合次序：本文分析。}
对本项目，一个逻辑完整的组合是：先由物理条件界定可识别的最简树和观测模型，再用 S14 组织损失权重或用 S16 的分裂坐标组织树搜索；如果具备可验证的正确似然与全域优化界，可进一步研究 S15 的置信集合；若任务是未来末端电压，则用 S17 在相应交换条件下校准预测区域。每次从一种输出转成另一种保证，都要补充新的统计条件。后验集中、搜索收敛、预测覆盖和真树保留，分别是四个需要单独验证的命题。
